\chapter{Delta Live Tables: Declarative Pipelines}
\index{Delta Live Tables}\index{declarative pipelines}\index{streaming table}\index{materialized view!DLT}%
\index{expectations}\index{data quality}\index{event log}\index{triggered pipeline}\index{continuous pipeline}%
\begin{objectives}
\begin{itemize}
    \item Explain DLT's declarative model: datasets, dependencies, and the automatically built DAG.
    \item Author streaming tables and materialized views in Python and SQL.
    \item Enforce data quality with expectations: warn, drop, and fail semantics.
    \item Operate pipelines: development versus production mode, triggered versus continuous, full refresh, and schema evolution.
    \item Query the DLT event log for quality metrics, lineage, and operational history.
    \item Decide deliberately between DLT and hand-written Spark plus Workflows.
\end{itemize}
\end{objectives}

\begin{crosscloud}
DLT is Databricks' \textbf{declarative pipeline} layer; the pattern---define datasets and quality,
let the platform handle execution---has analogs everywhere.

\vspace{2mm}
{\footnotesize
\begin{tabular}{@{}p{3.0cm}p{3.4cm}p{3.2cm}p{3.2cm}@{}}
\toprule
\textbf{Capability} & \textbf{Databricks (DLT)} & \textbf{AWS} & \textbf{Azure / GCP / Fabric} \\
\midrule
Declarative pipelines & DLT (Lakeflow Declarative Pipelines) & Glue workflows + DQ rules & ADF/Fabric pipelines + gates / Dataform \\
Quality rules & Expectations (warn/drop/fail) & Glue Data Quality (Deequ) & Great Expectations / dbt tests \\
Streaming dataset & Streaming table & Kinesis $\rightarrow$ Glue streaming & Eventstream / Dataflow / Fabric streams \\
Batch dataset & Materialized view & Redshift MV / Athena CTAS & BigQuery MV / Fabric MV \\
Ops telemetry & Event log (Delta table) & CloudWatch + Glue job metrics & ADF monitoring / Dataflow metrics \\
Lineage & Automatic, into UC & Glue lineage (partial) & Purview / Dataplex \\
\bottomrule
\end{tabular}}

\vspace{1mm}
\textbf{Snowflake}'s analog is Dynamic Tables with \texttt{EXPECT} clauses emerging; dbt
(cross-platform) shares the declarative-DAG philosophy with SQL models and tests.
\end{crosscloud}

\section{Week 10 Roadmap and Mental Model}
Chapter 9 squeezed performance out of hand-written Spark. Week 10 asks the architectural question: which pipelines should be hand-written at all? Delta Live Tables (DLT---being rebranded Lakeflow Declarative Pipelines) inverts control: you declare \emph{datasets} (streaming tables, materialized views) and their \emph{quality rules}; the framework builds the DAG, chooses batch or streaming execution, retries, checkpoints, and records every quality metric \cite{databricksDLT}.

The mental model: \textbf{DLT is a compiler from intent to pipeline}. Your code is a graph of dataset definitions; the runtime infers dependencies from table references, plans the execution order, and operates the result---including the parts hand-written pipelines get wrong under pressure: consistent checkpoints, quality telemetry, and lineage into Unity Catalog.

\begin{definitionbox}
\textbf{Delta Live Tables} is Databricks' declarative ETL framework. A \textbf{streaming table} ingests each source record exactly once (append processing); a \textbf{materialized view} maintains a query result, refreshed incrementally where possible. \textbf{Expectations} are per-dataset quality rules with warn, drop, or fail semantics.
\end{definitionbox}

\begin{keyterms}
\textbf{Expectation}: \texttt{@dlt.expect}, \texttt{expect\_or\_drop}, \texttt{expect\_or\_fail}. \textbf{Event log}: the per-pipeline Delta table of runs, quality metrics, and lineage. \textbf{Development mode}: fast-iterating pipeline mode that reuses compute and skips retries. \textbf{Full refresh}: rebuilding datasets from scratch. \textbf{Trigger}: available-now (triggered) versus continuous processing.
\end{keyterms}

\section{Monday: Datasets and the Declarative DAG}
\subsection{Streaming Tables and Materialized Views}
A streaming table wraps a streaming read---the append-only facts of your system (Auto Loader output, Kafka, CDC feeds). A materialized view wraps a query---DLT tracks its inputs and refreshes it incrementally when the plan allows. Dependencies are inferred: referencing \texttt{LIVE} names (or the table names) wires the DAG automatically.

\begin{lstlisting}[language=Python,caption={A three-dataset DLT pipeline in Python.},label={lst:ch10-pipeline}]
import dlt
from pyspark.sql import functions as F

@dlt.table(comment="Raw orders, exactly-once from the landing volume")
def bronze_orders():
    return (spark.readStream.format("cloudFiles")
            .option("cloudFiles.format", "csv")
            .option("header", True)
            .load("/Volumes/de_training/bronze/landing/olist/orders/"))

@dlt.table(comment="Typed, deduplicated orders")
def silver_orders():
    return (dlt.read_stream("bronze_orders")
            .withColumn("order_ts", F.to_timestamp("order_purchase_timestamp"))
            .withColumn("total_amount", F.col("payment_value").cast("decimal(12,2)"))
            .dropDuplicates(["order_id"]))

@dlt.table(comment="Daily revenue by region")
def gold_daily_revenue():
    return (dlt.read("silver_orders")
            .groupBy(F.to_date("order_ts").alias("order_date"), "region")
            .agg(F.sum("total_amount").alias("revenue")))
\end{lstlisting}

\subsection{Development versus Production Mode}
Development mode attaches to a long-lived cluster, skips retries, and reuses compute between runs---fast iteration. Production mode isolates each run, retries on failure, and terminates compute---operational safety. Develop in dev mode; deploy in production mode; never let a demo in dev mode stand in for a production run.

\section{Tuesday: Expectations---Quality as Code}
\subsection{Warn, Drop, Fail}
Expectations attach quality rules to datasets \cite{databricksDLTExpectations}: \texttt{expect} records violations (warn), \texttt{expect\_or\_drop} removes offending rows, \texttt{expect\_or\_fail} stops the update. All three emit metrics to the event log, so ``how much bad data did we see this week'' is a SQL query, not an investigation.

\begin{lstlisting}[language=Python,caption={The three expectation severities, deliberately assigned.},label={lst:ch10-expectations}]
@dlt.table
@dlt.expect("reasonable_amount", "total_amount < 1_000_000")   -- warn: log it
@dlt.expect_or_drop("valid_key", "order_id IS NOT NULL")       -- drop: unusable
@dlt.expect_or_fail("pii_gate", "customer_email IS NOT NULL")  -- fail: stop the line
def silver_orders_validated():
    return dlt.read_stream("bronze_orders")  # ... typed transform as above
\end{lstlisting}

\begin{pitfall}
\textbf{Everything is \texttt{expect}.} Warn-only expectations produce metrics nobody acts on---aspirational quality control. The discipline: every rule gets a severity decision on purpose. Keys and PII fail; business-plausibility rules warn with an owner; everything in between drops with quarantine visibility.
\end{pitfall}

\section{Wednesday: Operating Pipelines}
\subsection{Triggered versus Continuous}
A \textbf{triggered} pipeline processes available data and stops---batch economics for streaming sources. \textbf{Continuous} stays up for second-level latency at continuous cost. Most pipelines in this program are triggered on a schedule; Chapter 16 covers when continuous earns its DBUs.

\subsection{Full Refresh and Schema Evolution}
Incrementalization has limits: a logic change to a materialized view may require a \textbf{full refresh} (rebuild from source). Schema evolution of sources is handled per dataset (DLT tracks and adapts, or fails with a clear error depending on configuration). The operational habit: know per dataset whether tomorrow's change deploys cleanly or needs a refresh---and budget the refresh window.

\subsection{The Event Log}
Every pipeline publishes an event log (a Delta table in the pipeline's storage): run lifecycle, per-dataset quality metrics, lineage edges, and errors.

\begin{lstlisting}[language=SQL,caption={Quality metrics from the event log.},label={lst:ch10-eventlog}]
SELECT timestamp, details:flow_progress.data_quality.expectations
FROM event_log(TABLE(de_training.silver.silver_orders_validated))
WHERE event_type = 'flow_progress'
ORDER BY timestamp DESC
LIMIT 20;
\end{lstlisting}

\section{Thursday Hands-on Project: The Medallion as DLT}
Rebuild the Month 1 Bronze$\rightarrow$Silver$\rightarrow$Gold flow as one DLT pipeline with expectations, inject a poison record, and read the quality evidence from the event log.
\index{hands-on lab}\index{expectations!lab}

\begin{lstlisting}[language=Python,caption={Week 10 lab core: pipeline with poison-record drill.},label={lst:ch10-lab}]
# Pipeline notebook (selected datasets shown)
@dlt.table
@dlt.expect_or_drop("valid_order", "order_id IS NOT NULL")
@dlt.expect_or_fail("non_negative", "total_amount >= 0")
def silver_orders():
    return (dlt.read_stream("bronze_orders")
            .withColumn("total_amount", F.col("payment_value").cast("decimal(12,2)"))
            .dropDuplicates(["order_id"]))

# The drill: append one CSV with a NULL order_id and one negative amount.
# Expected: NULL-key row dropped (counted), negative amount fails the update.
# Then: fix the poison file, re-run, and confirm the pipeline recovers.
\end{lstlisting}

\subsection*{Deliverables}
\begin{itemize}
    \item The DLT pipeline (Python or SQL) reproducing the medallion flow with expectations on Silver.
    \item Event-log queries showing the dropped NULL-key record and the failed negative-amount update.
    \item A recovery note: what the pipeline did, what you fixed, how the re-run converged.
\end{itemize}

\subsection*{Acceptance Criteria}
\begin{itemize}
    \item The poisoned NULL-key row is dropped and \emph{counted} in quality metrics---not silently absent.
    \item The \texttt{expect\_or\_fail} violation stops the update; Gold is untouched by the failed run.
    \item After fixing the source, one triggered run converges Silver and Gold to the correct state with no manual table surgery.
    \item The pipeline runs in production mode with a schedule, not only interactively.
\end{itemize}

\subsection*{Stretch Goals}
\begin{itemize}
    \item Add a quarantine \emph{dataset} capturing dropped rows with reasons, mirroring Chapter 3's pattern.
    \item Convert \texttt{gold\_daily\_revenue} to a materialized view and document when it refreshes incrementally vs.\ fully.
    \item Deploy the pipeline twice (dev/prod targets) from one definition---a preview of Chapter 23.
\end{itemize}

\section{Friday Use Case: DLT or Hand-Written---Drawing the Line}
A logistics company asks you to classify three pipeline candidates: (1) 200-table nightly CDC replication, (2) a fraud feature pipeline with exotic stateful logic, (3) a compliance reporting flow with strict audit requirements.
\index{DLT!adoption}\index{architecture decision}

\subsection{Requirements and Assumptions}
The team knows Spark well and DLT not at all. The fraud pipeline uses \texttt{flatMapGroupsWithState} with custom timers. The compliance flow needs immutability evidence and sign-off gates.

\begin{figure}[H]
    \centering
    \resizebox{\textwidth}{!}{%
    \begin{tikzpicture}[
        cand/.style={rectangle, rounded corners, draw=primary, thick, fill=primary!7, align=center, minimum width=3.6cm, minimum height=1.0cm, font=\small\bfseries},
        verdict/.style={rectangle, rounded corners, draw=codegreen!60!black, thick, fill=green!8, align=center, minimum width=3.2cm, minimum height=0.9cm, font=\scriptsize\bfseries},
        why/.style={rectangle, rounded corners, draw=codegray, thick, fill=white, align=center, minimum width=3.6cm, minimum height=0.75cm, font=\scriptsize},
        flow/.style={-{Latex[length=2mm]}, draw=primary, thick}
    ]
        \node[cand] (cdc) at (0,0) {200-table CDC\\replication};
        \node[cand] (fraud) at (4.8,0) {Fraud features\\stateful logic};
        \node[cand] (comp) at (9.6,0) {Compliance\\reporting};
        \node[verdict] (v1) at (0,-1.5) {DLT + AUTO CDC};
        \node[verdict] (v2) at (4.8,-1.5) {Hand-written\\Spark + Workflows};
        \node[verdict] (v3) at (9.6,-1.5) {DLT + expectations\\+ event-log audit};
        \node[why] (w1) at (0,-2.8) {declared schemas,\\built-in SCD, scale};
        \node[why] (w2) at (4.8,-2.8) {custom state API\\outside DLT's model};
        \node[why] (w3) at (9.6,-2.8) {quality evidence is\\the deliverable};
        \draw[flow] (cdc) -- (v1); \draw[flow] (fraud) -- (v2); \draw[flow] (comp) -- (v3);
        \draw[flow, dashed] (v1) -- (w1); \draw[flow, dashed] (v2) -- (w2); \draw[flow, dashed] (v3) -- (w3);
    \end{tikzpicture}%
    }
    \caption{The adoption boundary: DLT where declaration wins, hand-written Spark where control wins.}
    \label{fig:ch10-boundary}
\end{figure}

\subsection{Architecture Decisions}
\textbf{CDC replication: DLT.} Two hundred tables of key-based upserts with SCD handling is exactly what \texttt{AUTO CDC} and expectations commoditize (Chapter 11 details the CDC mechanics). \textbf{Fraud features: hand-written.} Arbitrary stateful operators, custom timers, and side-effecting sinks live outside DLT's dataset model; forcing them in costs more than the DAG automation saves. \textbf{Compliance: DLT.} The event log plus expectations \emph{is} the audit evidence the requirement asks for.

\begin{pitfall}
\textbf{All-or-nothing adoption.} The platforms compose: a Workflows job (Chapter 12) can run a DLT pipeline task \emph{and} a hand-written Spark task. The boundary belongs inside the estate, per pipeline---not around it.
\end{pitfall}

\subsection{Risks and Mitigations}
\begin{table}[H]
\centering
\small
\begin{tabular}{@{}p{4.6cm}p{7.6cm}@{}}
\toprule
\textbf{Risk} & \textbf{Mitigation} \\
\midrule
Expectation drops quietly shrink a dataset & Drop-rate trend alert per dataset from the event log \\
Dev-mode pipeline promoted untested & Prod-mode deploy only via the bundle pipeline \\
Full refresh blows the weekend window & Refresh cost measured; calendar holds the slot \\
Framework upgrade changes behavior & Channel pinned; upgrade tested in dev target first \\
DLT hides logic from reviewers & Pipeline source is code in Git, reviewed like any change \\
\bottomrule
\end{tabular}
\caption{Week 10 scenario risks and mitigations.}
\label{tab:ch10-risks}
\end{table}

\section{Design Decision Guide}
\index{decision guide!Week 10}%
DLT decisions are about who operates the pipeline: you, or the framework.

\begin{table}[H]
\centering
\small
\begin{tabular}{@{}p{3.4cm}p{4.4cm}p{4.6cm}@{}}
\toprule
\textbf{Decision} & \textbf{Choose the first when} & \textbf{Choose the second when} \\
\midrule
DLT vs.\ hand-written Spark + Workflows & Standard ETL with quality rules & Exotic stateful logic or side effects \\
Streaming table vs.\ materialized view & Append-only facts, exactly-once ingest & Derived query results over inputs \\
\texttt{expect\_or\_drop} vs.\ \texttt{expect\_or\_fail} & Row is expendable, pipeline must continue & Bad data invalidates the whole batch \\
Triggered vs.\ continuous & Minutes-freshness, cost-sensitive & Measured sub-minute mandate \\
Development vs.\ production mode & Iterating & Anything scheduled or shared \\
Full refresh vs.\ incremental & Logic change invalidates stored state & Schema-compatible evolution \\
\bottomrule
\end{tabular}
\caption{Week 10 decision guide.}
\label{tab:ch10-decisions}
\end{table}

When torn, ask who will debug this at 03:00 and what evidence they will need. DLT's event log is often the deciding vote---operability is a feature of the platform, not of your diligence.

\section{Operational Guidance for Week 10}
\index{operational guidance!Week 10}%
\subsection{Performance and Reliability}
Track per-dataset processing time and quality metrics from the event log as first-class telemetry---the pipeline that degrades announces itself there weeks before users notice. Budget and calendar the full refreshes: know which logic changes force them and how long they take at current volume. Pin the pipeline channel deliberately and test upgrades in dev mode; a channel auto-upgrade before a demo is a rite of passage you can skip.

\subsection{Security and Governance Checklist}
\begin{itemize}
    \item Pipelines run as a service principal in production mode, never as a person.
    \item Expectation severities reviewed per rule: fail for keys/PII, drop with quarantine, warn with an owner.
    \item Event-log retention covers your audit window; export quality metrics for compliance.
    \item Pipeline target schemas are UC-governed like any table---declarative does not mean ungoverned.
\end{itemize}

\subsection{Troubleshooting Playbook}
\begin{table}[H]
\centering
\small
\begin{tabular}{@{}p{3.6cm}p{4.2cm}p{4.8cm}@{}}
\toprule
\textbf{Symptom} & \textbf{Likely cause} & \textbf{First action} \\
\midrule
Update fails on schema change & Source evolution incompatible with dataset & Read the error's dataset; evolve or full-refresh \\
Expectation storm after source change & Upstream contract broke & Event log quality trend; pause downstream publish \\
Pipeline ``stuck'' starting & Cluster quota or dev-mode reuse confusion & Check compute events; verify mode and policy \\
MV refreshed fully (slow) & Change invalidated incremental plan & Expected on logic change; schedule the window \\
\bottomrule
\end{tabular}
\caption{Week 10 troubleshooting quick reference.}
\label{tab:ch10-trouble}
\end{table}

\section{Interview Q\&A}
\subsection{Concepts and Trade-offs}
\begin{enumerate}
    \item \textbf{Q: What does DLT actually do that a notebook chain does not?}\\
    A: Dependency inference and DAG construction, execution-mode management (batch/streaming per dataset), checkpointing, retries, per-dataset quality metrics, an event log, and UC lineage---the operational layer you would otherwise hand-build.
    \item \textbf{Q: Streaming table versus materialized view?}\\
    A: A streaming table appends each source record exactly once (facts); a materialized view maintains a query result over its inputs (derived state). Streaming tables answer ``what arrived''; MVs answer ``what is true.''
    \item \textbf{Q: \texttt{expect\_or\_drop} versus \texttt{expect\_or\_fail}?}\\
    A: Drop when the row is expendable and the pipeline must continue (with quarantine visibility); fail when bad data invalidates the whole batch---keys, PII gates---and continuing would corrupt consumers.
    \item \textbf{Q: When does a DLT change require a full refresh?}\\
    A: When the change invalidates incremental state---logic changes to a materialized view whose stored results encode the old logic. Know it per dataset before you deploy.
    \item \textbf{Q: Why query the event log instead of trusting the UI?}\\
    A: The UI shows the last run; the event log is the durable, queryable record---trend analysis, audit evidence, and alerting all read it as a table.
    \item \textbf{Q: What telemetry would you lose by leaving DLT?}\\
    A: Per-dataset quality metrics, the event log's run history, and automatic lineage---you would rebuild all three before your first audit.
    \item \textbf{Q: How do expectations interact with schema evolution?}\\
    A: Expectations evaluate the row as it arrives at the dataset; a schema change that renames a column makes the expectation reference fail loudly---which is the point: quality rules pin the contract.
    \item \textbf{Q: Can DLT and plain Spark coexist in one pipeline estate?}\\
    A: That is the normal end state: DLT owns the standard flows, hand-written Spark owns the exotic ones, and Workflows orchestrates both as task types in the same jobs.
\end{enumerate}

\section{Further Reading}
The DLT documentation is the reference \cite{databricksDLT}; expectations and the event log are detailed in \cite{databricksDLTExpectations}. Auto Loader (the streaming source used throughout) is covered in \cite{databricksAutoLoader}, and the medallion rationale in \cite{databricksMedallion}. dbt practitioners should map the concepts across: datasets $\sim$ models, expectations $\sim$ tests, event log $\sim$ run results.

\section*{Key Takeaways}
\begin{itemize}
    \item DLT compiles declared datasets and quality rules into an operated pipeline---DAG, checkpoints, retries, metrics, lineage included.
    \item Streaming tables capture facts exactly once; materialized views maintain derived truth.
    \item Every expectation gets a deliberate severity: warn, drop (visible), or fail (protect downstream).
    \item The event log is your quality evidence base---query it, alert on it, audit with it.
    \item The DLT boundary is per pipeline: declaration wins for standard ETL; hand-written Spark wins for exotic control.
\end{itemize}

\section{Exercises}
\begin{enumerate}
    \item \textbf{(Conceptual)} Explain the difference between DLT's DAG inference and Airflow-style explicit dependencies; name one risk of each.
    \item \textbf{(Conceptual)} Why does \texttt{expect\_or\_fail} on a PII column protect \emph{downstream} tables, not just the dataset itself?
    \item \textbf{(Hands-on)} Complete the Thursday lab, including the poison drill and event-log queries.
    \item \textbf{(Hands-on)} Rewrite one dataset in SQL (DLT SQL) and compare readability with the Python version.
    \item \textbf{(Hands-on)} Trend this week's quality metrics from the event log and chart the drop rate per day.
    \item \textbf{(Design)} Classify five of your current (or hypothetical) pipelines on the DLT boundary; defend each call in one sentence.
    \item \textbf{(Operations)} Write the runbook for ``pipeline failed on \texttt{expect\_or\_fail} at 03:00'': diagnosis, fix, re-run, evidence.
    \item \textbf{(Architecture)} When would you graduate this pipeline's scheduling from the DLT trigger to a Workflows job (Chapter 12)?
\end{enumerate}
